\chapter{Clase 15}

\section{Deformaciones homogéneas}

Consideramos un sólido elástico lineal, homogéneo e isótropo, con constantes de Lamé $\lambda$ y $\mu$. Usamos la convención de suma de Einstein, escribimos los vectores como $\vec u$ y el tensor de deformación como $\mathbb U$, con componentes $U_{ij}$. En coordenadas cartesianas, $\nabla_i=\partial/\partial x_i$.

Recordemos las ecuaciones de equilibrio y la definición del tensor de deformación infinitesimal:
\[
f_i+\nabla_j\sigma_{ij}=0,
\qquad
U_{ij}=\frac{1}{2}(\nabla_i u_j+\nabla_j u_i).
\]
La ley de Hooke y su forma inversa son
\[
\sigma_{ij}=\lambda U_{kk}\delta_{ij}+2\mu U_{ij},
\]
\[
U_{ij}=\frac{\sigma_{kk}}{9K}\delta_{ij}
+\frac{1}{2\mu}\left(\sigma_{ij}-\frac{\sigma_{kk}}{3}\delta_{ij}\right),
\qquad K=\lambda+\frac{2}{3}\mu.
\]
Aquí, $U_{kk}=\mathrm{Tr}\,\mathbb U$ y $\sigma_{kk}$ es la traza del tensor de esfuerzos. Los esfuerzos normales positivos corresponden a tracción y los negativos a compresión.

Una deformación es \emph{homogénea} cuando $U_{ij}$ es independiente de la posición:
\[
\nabla_k U_{ij}=0.
\]
Como los módulos elásticos son constantes, la ley de Hooke implica que $\nabla_k\sigma_{ij}=0$. La ecuación de equilibrio exige entonces $f_i=0$: bajo estas hipótesis, una deformación homogénea en equilibrio no admite fuerzas de cuerpo no nulas.

Las condiciones de frontera se imponen sobre el vector de tracción:
\[
t_i=\sigma_{ij}n_j,
\]
donde $\hat n$ es la normal exterior. En una interfaz sin fuerzas superficiales adicionales, la tracción es continua; usando una misma normal a ambos lados, esta condición se escribe $[\sigma_{ij}n_j]=0$, donde los corchetes indican el salto entre ambos medios.

\subsection{Compresión hidrostática}

Una presión uniforme $p>0$ sobre toda la superficie produce la tracción $t_i=-pn_i$. El estado homogéneo que satisface esta condición es
\[
\sigma_{ij}=-p\delta_{ij},
\qquad \sigma_{kk}=-3p.
\]
Al sustituir en la ley de Hooke inversa, la parte desviadora del esfuerzo se anula:
\[
\begin{aligned}
U_{ij}
&=-\frac{p}{3K}\delta_{ij}
+\frac{1}{2\mu}\left(-p\delta_{ij}+p\delta_{ij}\right)\\
&=\boxed{-\frac{p}{3K}\delta_{ij}}.
\end{aligned}
\]
Cada deformación normal vale $-p/(3K)$, pero la traza suma las tres direcciones:
\[
U_{xx}=U_{yy}=U_{zz}=-\frac{p}{3K},
\qquad
\frac{\Delta V}{V}=U_{ii}=-\frac{p}{K}.
\]
El cambio de volumen se expresa a primer orden en la deformación. Un campo de desplazamiento que genera este tensor es
\[
u_i=-\frac{p}{3K}x_i+C_i.
\]
La parte proporcional a $x_i$ representa la compresión; $C_i$ representa una traslación rígida. Se omite una posible rotación rígida infinitesimal, que tampoco modifica $U_{ij}$.

\subsection{Estiramiento uniforme}

Consideremos una varilla larga orientada en el eje $z$, sometida a tracción uniforme en sus extremos y sin fuerzas sobre las caras laterales. Para una normal lateral $\hat n=(n_x,n_y,0)$,
\[
t_i=\sigma_{ix}n_x+\sigma_{iy}n_y=0.
\]
Como las normales laterales abarcan las direcciones del plano $xy$, el estado homogéneo debe satisfacer $\sigma_{ix}=\sigma_{iy}=0$. Por la simetría del tensor de esfuerzos, también se anulan $\sigma_{xi}$ y $\sigma_{yi}$. La única componente no nula es
\[
\sigma_{zz}=p>0,
\qquad \sigma_{kk}=p.
\]
Aquí $p$ denota la magnitud del esfuerzo de tracción, no una presión compresiva. Las tracciones en los extremos son $p\hat e_z$ y $-p\hat e_z$, de acuerdo con sus normales exteriores.

La ley de Hooke inversa da
\[
U_{ij}=\frac{p}{9K}\delta_{ij}
+\frac{1}{2\mu}\left(\sigma_{ij}-\frac{p}{3}\delta_{ij}\right).
\]
Por tanto,
\[
\begin{aligned}
U_{xx}=U_{yy}
&=\left(\frac{1}{9K}-\frac{1}{6\mu}\right)p
=-\frac{3K-2\mu}{18K\mu}p,\\
U_{zz}
&=\left(\frac{1}{9K}+\frac{1}{3\mu}\right)p
=\frac{3K+\mu}{9K\mu}p.
\end{aligned}
\]
Definimos el \emph{módulo de Young} $E$ y el \emph{coeficiente de Poisson} $\nu$ mediante
\[
\boxed{E=\frac{9K\mu}{3K+\mu}},
\qquad
\boxed{\nu=\frac{3K-2\mu}{2(3K+\mu)}}.
\]
Con estas definiciones, las deformaciones se escriben de forma sencilla:
\[
U_{zz}=\frac{p}{E},
\qquad U_{xx}=U_{yy}=-\nu\frac{p}{E},
\qquad \nu=-\frac{U_{xx}}{U_{zz}}.
\]
Para $\nu>0$, la varilla se alarga longitudinalmente y se contrae transversalmente.

\textbf{Nota.} El módulo de Young de los metales suele ser del orden de $100\,\mathrm{GPa}$ y su coeficiente de Poisson suele estar cerca de $0.3$. La estabilidad elástica, $K>0$ y $\mu>0$, exige $-1<\nu<1/2$; los extremos corresponden a límites, no a valores admisibles con ambos módulos finitos y positivos.

El cambio relativo de volumen y la energía elástica por unidad de volumen de referencia son
\[
\frac{\Delta V}{V}=U_{ii}
=(1-2\nu)\frac{p}{E}=\frac{p}{3K},
\qquad
F=\frac{1}{2}\sigma_{ij}U_{ij}
=\frac{p^2}{2E}.
\]
Tomando como referencia un campo sin traslaciones ni rotaciones rígidas, los desplazamientos son
\[
u_x=-\nu\frac{p}{E}x,
\qquad u_y=-\nu\frac{p}{E}y,
\qquad u_z=\frac{p}{E}z.
\]

\subsection{Relaciones entre las constantes elásticas}

Las constantes de Lamé y el módulo volumétrico pueden expresarse en términos de $E$ y $\nu$:
\[
\lambda=\frac{E\nu}{(1+\nu)(1-2\nu)},
\qquad
\mu=\frac{E}{2(1+\nu)},
\qquad
K=\frac{E}{3(1-2\nu)}.
\]
Así, las formas directa e inversa de la ley de Hooke quedan
\[
\sigma_{ij}=\frac{E}{1+\nu}
\left(U_{ij}+\frac{\nu}{1-2\nu}U_{kk}\delta_{ij}\right),
\]
\[
U_{ij}=\frac{1}{E}
\left[(1+\nu)\sigma_{ij}-\nu\sigma_{kk}\delta_{ij}\right].
\]

\subsection{Compresión con restricción lateral}

Consideremos una varilla comprimida a lo largo del eje $z$, cuyos lados están restringidos para impedir el desplazamiento transversal. A diferencia del caso con caras laterales libres, aquí aparecen esfuerzos laterales que mantienen
\[
U_{xx}=U_{yy}=0,
\qquad U_{ij}=0\quad\text{si }i\ne j.
\]
La única deformación no nula es $U_{zz}$, de modo que $U_{kk}=U_{zz}$. La ley de Hooke da
\[
\begin{aligned}
\sigma_{xx}=\sigma_{yy}
&=\lambda U_{zz}
=\frac{E\nu}{(1+\nu)(1-2\nu)}U_{zz},\\
\sigma_{zz}
&=(\lambda+2\mu)U_{zz}
=\frac{E(1-\nu)}{(1+\nu)(1-2\nu)}U_{zz}.
\end{aligned}
\]
Si el esfuerzo axial aplicado es $\sigma_{zz}=-p$, con $p>0$,
\[
\boxed{U_{zz}=-\frac{p(1+\nu)(1-2\nu)}{E(1-\nu)}}.
\]
Integrando y fijando $u_z(0)=0$, obtenemos
\[
u_x=u_y=0,
\qquad
u_z=-\frac{p(1+\nu)(1-2\nu)}{E(1-\nu)}z.
\]
Los esfuerzos laterales asociados a la restricción son
\[
\boxed{\sigma_{xx}=\sigma_{yy}=-\frac{p\nu}{1-\nu}}.
\]
Para $\nu>0$, también son compresivos: impiden la expansión transversal que ocurriría si los lados estuvieran libres.

\subsection{Cizalla uniforme}

Consideremos un bloque sometido a un estado de corte uniforme cuyas únicas componentes de esfuerzo no nulas son
\[
\sigma_{xy}=\sigma_{yx}=P.
\]
Como $\sigma_{kk}=0$, la ley de Hooke implica
\[
U_{xy}=U_{yx}=\frac{P}{2\mu},
\qquad
\mathbb U=\begin{pmatrix}
0&P/(2\mu)&0\\
P/(2\mu)&0&0\\
0&0&0
\end{pmatrix}.
\]
La definición del tensor de deformación exige
\[
\frac{1}{2}(\nabla_x u_y+\nabla_y u_x)=\frac{P}{2\mu}.
\]
Una elección que satisface esta relación es el campo de cizalla simple
\[
u_x=\frac{P}{\mu}y,
\qquad u_y=u_z=0.
\]
El ángulo de corte infinitesimal y la rotación local alrededor del eje $z$ son, respectivamente,
\[
\gamma=2U_{xy}=\frac{P}{\mu},
\qquad
\phi=\frac{1}{2}(\nabla_x u_y-\nabla_y u_x)
=-\frac{P}{2\mu}.
\]
Así, este campo combina deformación de corte y rotación local; el signo negativo indica el sentido de giro respecto de $+\hat e_z$. Como $U_{ii}=0$, no hay cambio de volumen a primer orden.

\section{Reconstrucción del tensor de esfuerzos}

Consideremos el campo de desplazamiento cartesiano
\[
u_1=kx_2x_3,
\qquad u_2=kx_1x_3,
\qquad u_3=kx_1x_2,
\]
donde $k$ es constante y los gradientes de desplazamiento son pequeños. Usamos $x_i$ para las coordenadas cartesianas, de manera consistente con la notación anterior.

Calculamos las seis componentes independientes de $\mathbb U$:
\[
\begin{aligned}
U_{11}&=\nabla_1 u_1=0,\\
U_{22}&=\nabla_2 u_2=0,\\
U_{33}&=\nabla_3 u_3=0,\\
U_{12}=U_{21}
&=\frac{1}{2}(\nabla_1 u_2+\nabla_2 u_1)
=\frac{1}{2}(kx_3+kx_3)=kx_3,\\
U_{13}=U_{31}
&=\frac{1}{2}(\nabla_1 u_3+\nabla_3 u_1)
=\frac{1}{2}(kx_2+kx_2)=kx_2,\\
U_{23}=U_{32}
&=\frac{1}{2}(\nabla_2 u_3+\nabla_3 u_2)
=\frac{1}{2}(kx_1+kx_1)=kx_1.
\end{aligned}
\]
Por tanto,
\[
\mathbb U=k\begin{pmatrix}
0&x_3&x_2\\
x_3&0&x_1\\
x_2&x_1&0
\end{pmatrix},
\qquad \mathrm{Tr}\,\mathbb U=U_{ii}=0.
\]
Al sustituir en la ley de Hooke, el término volumétrico desaparece:
\[
\sigma_{ij}=\frac{E}{1+\nu}
\left(U_{ij}+\frac{\nu}{1-2\nu}U_{kk}\delta_{ij}\right)
=\frac{E}{1+\nu}U_{ij}=2\mu U_{ij}.
\]
Las componentes del tensor de esfuerzos se organizan entonces en la matriz
\[
\boxed{\{\sigma_{ij}\}=\frac{kE}{1+\nu}
\begin{pmatrix}
0&x_3&x_2\\
x_3&0&x_1\\
x_2&x_1&0
\end{pmatrix}}.
\]
Este ejemplo no es una deformación homogénea: $U_{ij}$ varía con la posición, aunque su traza sea nula.